\section{An actual whole-cluster lower bound for the quotient cost}
\label{sec:cluster-quotient-cost}

The single-jet estimate in MA-PR-012 can be replaced by
a functional on an entire cluster.  The functional has
polynomial norm in the original full block algebra,
irrespective of the separations inside that cluster.
Its composite with the canonical Hilbert return measures
an actual weighted exponential-polynomial norm.  This
gives a conditional obstruction to polynomial costs even
for clusters of distinct simple zeros.  No occurrence
of such clusters for the fixed arithmetic function is
asserted.

Retain the fixed data of (MA-PR.1), the original
MT-M18 rectangles and quotient norms, and the
no-proper-truncation condition (MA-PR.100).  Thus
\[
 \begin{gathered}
 e_-=a+u,\qquad e_+=a+v,\qquad p=\frac{s+1}{2},
 \qquad \nu_n=2+|n|,\\
 \mathcal C_n
   =\|\iota_n\mathscr R_n\|_{\mathcal L(Q_H,J_n)} .
 \end{gathered}
 \tag{MA-PR.175}
\]
The inclusion \(\iota_n\) is the identity at every
fully retained primary factor and zero at every
discarded factor.  The source norm in (MA-PR.175)
is the norm on the whole \(Q_H\).  No generic edge
distance, simplicity, or separation between zeros
is assumed below.

\paragraph{MA-PR-013 (cluster depth in the actual quotient norm).}
Let \(\rho_1,\ldots,\rho_k\), \(k\ge1\), be a multiset
of fully retained zeros in one original block
\(\Omega_n\).  A zero may be repeated no more than
its actual multiplicity.  Suppose that, for some
\(0<d\le1\) and \(T\in\mathbb R\),
\[
       |\rho_j-(e_+-d+iT)|\le d/2
                      \qquad(1\le j\le k).
 \tag{MA-PR.176}
\]
There are constants \(c_s>0\) and \(C,D>0\),
depending only on the fixed arithmetic function,
Hilbert parameters and selected block system, for which
\[
 \mathcal C_n\ge
       c_s C^{-1}\nu_n^{-D}
                   3^{-k}k^{-s/2}d^{\,p-k}.
 \tag{MA-PR.177}
\]
In particular, after changing the fixed constants,
\[
                  \mathcal C_n\ge
                         c\,\nu_n^{-D'}d^{\,p-k}.
 \tag{MA-PR.178}
\]
Both estimates hold with the disk in (MA-PR.176)
replaced by the disk centered at \(e_-+d+iT\)
of radius \(d/2\).  Constants are independent of
the cluster size, its internal separations, and
its depth \(d\).

Here is a consequence in which membership in a
single block need not be assumed in advance.
Suppose there is a sequence of multisets of actual
zeros in the disks of (MA-PR.176), with sizes \(k_j\),
centers \(T_j\), and depths \(d_j\), such that
\[
 \begin{gathered}
 |T_j|\longrightarrow\infty,\qquad
 d_j\longrightarrow0,\qquad
 \frac{(k_j-p)\log(1/d_j)}
                     {\log(2+|T_j|)}
                         \longrightarrow+\infty .
 \end{gathered}
 \tag{MA-PR.179}
\]
The same assertion is allowed at the left edge.
After a finite prefix, every such multiset lies
in one original block, consists of strictly
interior fully retained zeros, and forces failure
of every polynomial bound on \(\mathcal C_n\).
Consequently
\[
 \{g\in\mathcal D_{[1,2]}:
                    \sup_n c_n(g)=\infty\}
       \quad\hbox{is a dense }G_\delta
                       \hbox{ in }\mathcal D_{[1,2]} .
 \tag{MA-PR.180}
\]
Every test in that set has \(g(0)=0\) and has no
continuous equivariant return of (MA-PR.51).
This conclusion concerns the actual quotient-return
problem.  It does not identify it with the distinct
full common-space matrix problem of MA-MO-033.

\begin{proof}
\emph{1. A cluster functional on the original full block.}
Set \(P(z)=\prod_{j=1}^k(z-\rho_j)\), with all the
chosen repetitions retained.  For \(F\in A(\Omega_n)\)
define
\[
       \Delta_P[F]
          =\frac1{2\pi i}
                   \int_{\Gamma_n}\frac{F(z)}{P(z)}\,dz .
 \tag{MA-PR.181}
\]
If \(F\in I_n\), then it vanishes at each chosen
node to at least its chosen multiplicity.  Thus
\(F/P\) is holomorphic throughout the rectangle,
and (MA-PR.181) is zero.  The functional is therefore
well defined on the original \(J_n=A(\Omega_n)/I_n\).

On either vertical side of \(\Gamma_n\), every
factor \(|z-\rho_j|\) is at least one, since the
arithmetic zeros lie in \(0\le\operatorname{Re}z\le1\)
and the sides are at \(-1\) and \(2\).  On a
horizontal side at height \(h_l\), \(l=n,n+1\),
\[
 \begin{split}
 |P(z)|
 &\ge\prod_{j=1}^k
               \min\{1,|h_l-\operatorname{Im}\rho_j|\}\\
 &\ge\prod_{\rho:\,\Xi(\rho)=0}
               \min\{1,|h_l-\operatorname{Im}\rho|\}^{m_\rho}
       \ge c\,\nu_n^{-D}.
 \end{split}
 \tag{MA-PR.182}
\]
The middle inequality holds because omitted factors
are at most one and no chosen multiplicity exceeds
the actual multiplicity.  The final inequality is
exactly (M18.7), with the two adjacent height indices
comparable to \(n\).  The uniformly bounded perimeter
of \(\Gamma_n\) now gives
\(\|\Delta_P\|_{J_n^*}\le C\nu_n^D\), by taking the
infimum over all representatives \(F\) of a class.
This estimate has no inverse difference between
two cluster nodes.

\emph{2. Identify its composite on the whole Hilbert quotient.}
Define the entire exponential polynomial
\[
          K_P(x)=\frac1{2\pi i}
                    \int_{\Gamma_n}\frac{e^{zx}}{P(z)}\,dz .
 \tag{MA-PR.183}
\]
Its residue expansion uses derivatives at the chosen
zeros up to their chosen multiplicities minus one.
The identity of the canonical inclusion on those
full factors therefore implies
\[
 \Delta_P\iota_n\mathscr R_n(\pi h)
       =\int_{\mathbb R}h(x)e^{-ax}K_P(x)\,dx .
 \tag{MA-PR.184}
\]
One can first verify this for \(h\in\mathcal E\):
its Mellin transform is entire, and integration
over the fixed compact contour may be interchanged
with its absolutely convergent Mellin integral.
The residue formula gives exactly the chosen
finite linear combination of (MA-PR.46).
Every such jet is continuous on \(H\), since its
factor is fully retained.  Density of
\(\mathcal E\) in \(H\) then proves (MA-PR.184)
for every \(h\in H\).  In particular no Mellin
evaluation of a general \(h\in H\) on the outer
contour has been assumed.

The functional in (MA-PR.184) annihilates \(J_H\).
The dual of the Hilbert quotient is isometric to
that annihilator in \(H^*\), so its exact norm is
\[
 \|\Delta_P\iota_n\mathscr R_n\|_{Q_H^*}^{\,2}
     =\int_{\mathbb R}
          \frac{e^{-2ax}|K_P(x)|^2}
                {\max(e^{2ux},e^{2vx})(1+x^2)^{s/2}}\,dx .
 \tag{MA-PR.185}
\]
This equality uses the whole quotient norm, rather
than a norm imposed separately on the algebraic
coordinates \(B_n\).

\emph{3. A separation-free estimate of the positive half-line.}
Under (MA-PR.176), write
\(d_j=e_+-\operatorname{Re}\rho_j\) and
\(\epsilon_j=\operatorname{Im}\rho_j-T\).  Then
\(d/2\le d_j\le3d/2\) and \(|\epsilon_j|\le d/2\).
For \(x\ge0\) put
\[
 \begin{split}
 \kappa_P(x)&=e^{-(e_++iT)x}K_P(x),\\
 \kappa_P&=
       \mathop{*}_{j=1}^k
       \bigl(e^{-(d_j-i\epsilon_j)x}
                         {\bf1}_{[0,\infty)}(x)\bigr).
 \end{split}
 \tag{MA-PR.186}
\]
For completeness, the Laplace transform of
\(K_P|_{[0,\infty)}\), to the right of all its
exponents, is \(1/P(z)\).  To justify the initial
interchange on the whole outer contour, first take
\(\operatorname{Re}z>2\).  Integrating
\(e^{-zx}e^{wx}\) first in (MA-PR.183), the
resulting sum of residues of
\(1/((z-w)P(w))\) is \(1/P(z)\); holomorphy
continues this identity to the right of all exponents.
The convolution
of the \(k\) exponentials has the same Laplace
transform.  Uniqueness of the Fourier transform
after an exponentially decaying twist proves the
identity, including repeated nodes.  Each factor
in (MA-PR.186) is in \(L^1\cap L^2\); its
convolution also is.  With the unnormalized Fourier
transform \(\widehat f(\xi)=\int e^{-i\xi x}f(x)\,dx\),
\[
       \widehat\kappa_P(\xi)
            =\prod_{j=1}^k
                     \frac1{d_j+i(\xi-\epsilon_j)} .
 \tag{MA-PR.187}
\]
For \(|\xi|\le d\), each denominator has modulus
at most \(3d\).  Plancherel therefore gives
\[
     \|\kappa_P\|_2^2
        \ge\frac1\pi\,9^{-k}d^{\,1-2k}.
 \tag{MA-PR.188}
\]
This step handles arbitrary phases within the
disk.  No common-phase estimate of an oscillatory
simplex integral is needed.

The convolution simplex also gives the pointwise
upper bound
\[
       |\kappa_P(x)|
          \le e^{-dx/2}\frac{x^{k-1}}{(k-1)!}
                         \qquad(x\ge0).
 \tag{MA-PR.189}
\]
Let \(\Gamma(q,b)=\int_b^\infty t^{q-1}e^{-t}\,dt\).
The \(L^2\) mass after \(16k/d\) is consequently at most
\[
 \begin{split}
 \int_{16k/d}^\infty|\kappa_P(x)|^2\,dx
 &\le d^{\,1-2k}
                  \frac{\Gamma(2k-1,16k)}{((k-1)!)^2}\\
 &\le d^{\,1-2k}e^{-8k}
             2^{2k-1}\frac{(2k-2)!}{((k-1)!)^2}\\
 &\le \frac18(16e^{-8})^k d^{\,1-2k}
 \le \frac18\,9^{-k}d^{\,1-2k}.
 \end{split}
 \tag{MA-PR.190}
\]
The second inequality uses \(e^{-t}\le
e^{-8k}e^{-t/2}\) for \(t\ge16k\); the third uses
\(\binom{2k-2}{k-1}\le4^{k-1}\).
The last inequality follows from \(e^8>144\).
Combining (MA-PR.188)--(MA-PR.190), and using
\(\pi<4\), leaves at least
\(9^{-k}d^{1-2k}/8\) of the squared norm in
\([0,16k/d]\).

On \(x\ge0\) the maximum in the weight is
\(e^{2vx}\), so the positive half of
(MA-PR.185) is exactly the weighted squared
norm of \(\kappa_P\).  Since \(16k/d\ge1\),
\[
 \begin{split}
 \int_0^\infty
       |\kappa_P(x)|^2(1+x^2)^{-s/2}\,dx
 &\ge \frac{2^{-s/2}}{8\cdot16^s}\,
            9^{-k}k^{-s}d^{\,s+1-2k},\\
 \|\Delta_P\iota_n\mathscr R_n\|
 &\ge c_s\,3^{-k}k^{-s/2}d^{\,p-k}.
 \end{split}
 \tag{MA-PR.191}
\]
Together with the actual \(J_n^*\) bound following
(MA-PR.182), this proves (MA-PR.177).

The ordinary arithmetic block count (M18.2) gives
\(k\le C_0\log\nu_n\).  Thus \(3^{-k}\) is
bounded below by a fixed negative power of
\(\nu_n\), and, for each fixed \(s\), so is
\(k^{-s/2}\).  Absorbing these factors proves
(MA-PR.178) with constants independent of \(k\).

\emph{4. The left edge.}
If the cluster is centered at \(e_-+d+iT\), use
the negative half of (MA-PR.185), put \(x=-y\),
and remove the exponent \(e^{-(e_-+iT)y}\).
Changing the residue variable from \(z\) to \(-z\)
gives
\[
 e^{(e_-+iT)y}K_P(-y)
   =(-1)^{k-1}
      \mathop{*}_{j=1}^k
       \bigl(e^{-(\rho_j-e_--iT)y}
                         {\bf1}_{[0,\infty)}(y)\bigr).
 \tag{MA-PR.192}
\]
The real decay rates again lie between \(d/2\)
and \(3d/2\), and the imaginary deviations have
absolute value at most \(d/2\).  On the negative
half-line the maximum in the weight is
\(e^{2ux}\).  Hence it cancels the removed
exponent exactly as before.  All the estimates
(MA-PR.187)--(MA-PR.191) apply unchanged in
absolute value, proving the left-edge assertion.

\emph{5. An arbitrarily small bad cluster cannot cross a selected boundary.}
Consider (MA-PR.179), and delete a finite prefix
so that \(d_j<\min(1/4,2(v-u)/3)\).  The entire
disk then lies strictly between \(e_-\) and
\(e_+\).  Every zero in it is retained in full
by the actual Hilbert jet classification.
The chosen zero ordinates occupy an interval
of length at most \(d_j\).

If they do not all belong to one block, a
selected height \(h_l\) lies strictly between
two of their ordinates.  No ordinate equals
\(h_l\).  Every one of the \(k_j\) chosen
ordinates then has distance at most \(d_j\)
from \(h_l\), while \(2+|l|\) is comparable
to \(2+|T_j|\).  Applying (M18.7) to this
subset of factors gives
\[
       c(2+|T_j|)^{-D}
             \le d_j^{\,k_j},
 \qquad
       k_j\log(1/d_j)
             \le D\log(2+|T_j|)+O(1).
 \tag{MA-PR.193}
\]
This contradicts (MA-PR.179), since
\(k_j\ge k_j-p\) and \(p>0\).
Thus eventually every cluster lies in one
original block \(\Omega_{n_j}\), with
\(\nu_{n_j}\) comparable to \(2+|T_j|\).
The already proved estimate now gives
\[
 \frac{\log\mathcal C_{n_j}}
                   {\log(2+|T_j|)}
 \ge
 \frac{(k_j-p)\log(1/d_j)}
                   {\log(2+|T_j|)}
                  -D'+o(1)
       \longrightarrow+\infty .
 \tag{MA-PR.194}
\]
No polynomial cost bound is possible.  The
precisely quantified fixed-support Baire
alternative (MA-PR.104) proves (MA-PR.180).
The same boundary-crossing argument applies
to left-edge clusters.
\end{proof}

\paragraph{Exact scope and source placement.}
The implication just proved is conditional on
the actual cluster geometry (MA-PR.179).
Neither that geometry nor superpolynomial
canonical costs are asserted for the fixed
arithmetic zero divisor.  The ordinary local
zero count alone does not imply this geometry.
Nor does the theorem assert that the sufficient
potential of MA-PR-012 is necessary.

The only zero-count input in the lower bound is
the established block count and the original
selected-height product (M18.2), (M18.7).
The actual quotient and retained-jet interfaces
are (MA-PR.1), (MA-PR.45)--(MA-PR.47), and
the canonical no-proper-truncation inclusion
(MA-PR.100)--(MA-PR.106).  In particular the
argument has not changed block norms, chosen
separated scalar interpolation coordinates, or
deleted an oblique projection.

Repeated roots in \(P\) represent their actual
finite-length primary factors.  The residue
functional respects those full Taylor classes,
but its Hilbert dual norm contains additional
analytic information.  The finite-length,
Grothendieck-group and nuclear-trace scope is
the one already specified in (MA-PR.80) and
MA-PR-009--011; those algebraic descriptions
do not supply a polynomial norm estimate.
No fresh reading of the inherited cote~1
locators or new outside arithmetic premise
is claimed here.
